\chapter{Stato Attuale}

\section{Sistemi ERP}
I sistemi \textit{ERP} (\textit{Enterprise Resource Planning}) rappresentano un elemento essenziale del sistema gestionale aziendale \cite{erpMigration},
consentendo l'integrazione dei processi di \textit{business} principali di un'azienda come: vendite, acquisti, gestione inventario, contabilità e personale.

\subsection{Evoluzione dei Sistemi Legacy al Cloud}
Nel corso degli anni, i sistemi \textit{ERP} si sono evoluti partendo da architetture \textit{on-premise} \textit{legacy}, spesso caratterizzate da un elevato
grado di personalizzazione (\textit{customization}, in Infor LN definiti come \texttt{tx}) e importanti costi di manutenzione.
Tuttavia, oggi si sta assistendo ad una transizione verso nuove soluzioni \textit{Cloud}.

% \subsection{Architetture Multitenant}
Di fatto, il modello \textit{Cloud Multitenant} oggi rappresenta l'architettura degli \textit{ERP} moderni \cite{multitenantCloud}.
Ove, una singola istanza del software e infrastruttura servono più clienti (definiti come \textit{tenant}).
Questo approccio offre scalabilità, aggiornamenti continui e ottimizzazione dei costi, ma anche grandi sfide nell'ambito
della gestione e segregazione dei dati, nonché della sicurezza, richiedendo che ogni processo di migrazione operi
secondo specifici vincoli di integrità \cite{dataMigrationStrategies}.

\section{Infor LN e CloudSuite Enterprise}
Infor LN, nato con il nome di Baan \textit{ERP}, è uno dei principali \textit{ERP} nel settore manifatturiero industriale \cite{inforLN2025} a livello internazionale.
Viene integrato all'interno della Infor \textit{CloudSuite Enterprise}, il quale offre funzionalità per la gestione della \textit{supply chain}
e produzione \cite{inforCloudsuite}.

L'interfaccia utente Web (nota come \texttt{LNUI}) permette agli utenti finali di navigare tra le diverse sessioni
tramite \textit{menu} contestuali o ricerca diretta (Figura~\ref{fig:ln_ui}).

\begin{figure}[H]
    \centering
    \includegraphics[width=0.5\textwidth]{Lista_Sessioni_LNUI}
    \caption{Esempio di interfaccia Web Infor LN (LNUI) con elenco sessioni del Migration Tool.}
    \label{fig:ln_ui}
\end{figure}

% \subsection{Architettura e Componenti di Infor LN}
Infor LN è basato su un'architettura che separa il \textit{framework} tecnologico (gli Infor Tools) dalla logica di \textit{business}.
Ogni applicazione è composta da diversi oggetti:
\begin{itemize}
    \item \textbf{Tabelle (\textit{Tables}):} I dati fisici, definiti tramite \textit{Data Dictionary}, con domini standard o \textit{custom}.
    \item \textbf{Sessioni (\textit{Sessions}):} Le interfacce utente (da qui definite come Sessioni) attraverso cui gli operatori accedono, visualizzano e interagiscono con i dati.
    \item \textbf{DAL (\textit{Data Access Layer}):} L'elemento che si occupa di garantire l'integrità dei dati e delle tabelle, indipendentemente dalle Sessioni \cite{inforDalGuide}.
    \item \textbf{Domini (\textit{Domains}):} Definiscono i tipi di dati, con regole di validazione e formati, che possono essere standard o personalizzati (\textit{custom}).
    \item \textbf{Etichette (\textit{Labels}):} Utilizzate per la localizzazione e la visualizzazione dei campi nelle Sessioni, con supporto multilingue (es.
    02 corrisponde alla lingua inglese, e 06 alla lingua italiano).
\end{itemize}

\subsection{Linguaggio Baan C: 3GL e 4GL}
Per lo sviluppo e personalizzazione degli \textit{script} e librerie, Infor LN utilizza un linguaggio proprietario, di nome \textbf{Baan C} \cite{infor3GL4GL}.

Esso si divide a sua volta in:
\begin{itemize}
    \item \textbf{4GL (\textit{Fourth-Generation Language}):} La versione più recente e moderna.
    Infor LN Tools genera automaticamente operazioni standard di Sessione e database (CRUD). In 4GL non è necessario
    scrivere codice per gestire il salvataggio o il caricamento di un record, ma ci si concentra sulla gestione
    degli \textit{eventi} legati ai componenti della Sessione o del database (tramite \textit{DAL} \cite{inforDal2Guide}).
    \item \textbf{3GL (\textit{Third-Generation Language}):} Precedente a 4GL, e totalmente integrato con esso, viene utilizzato
    per la gestione di logiche di \textit{business} complesse.
    Permette un approccio più procedurale e a basso livello per implementare funzioni matematiche, algoritmi personalizzati,
    \textit{DLL} (\textit{Dynamic Link Libraries}) proprietarie, funzioni e routine di manipolazione stringhe o file, nonchè di estensioni (\textit{extensions})
    che rappresentano in alcuni casi, l'unico modo di interagire e personalizzare tabelle e procedure standard.
\end{itemize}

\subsection{Infor Exchange e Strumenti di Import/Export}
Per lo scambio di informazioni, Infor fornisce un modulo denominato \textit{Exchange} \cite{inforExchangeGuide}.
Si tratta di una piattaforma per l'importazione e l'esportazione di file (es.
CSV, XML, testo a larghezza fissa)
mappando righe e colonne sui \textit{Datafiles} e sui campi del database di LN .

Spesso sono necessarie delle personalizzazioni delle procedure di migrazione tramite Baan C, che quindi risulta essenziale quando le regole di
mappatura superano la complessità gestibile nativamente dall'\textit{Exchange}.

\section{Migrazione Dati nei Sistemi ERP}
\subsection{Sfide e Problematiche Comuni}
La migrazione da un sistema \textit{legacy} a un nuovo \textit{ERP} \textit{cloud} richiede un'attenta fase di analisi per evitare perdite o
disallineamenti di informazioni storiche e \textit{master data}.

Attualmente, molte migrazioni sono gestite manualmente, con l'integrazioni di \textit{script} scritti
\textit{ad hoc}, con un alto livello di complessità e rischio d'errore.
Per questo, nasce l'idea del Tool di Migrazione (\textit{Migration Tool}) oggetto di questa tesi.

%\subsection{Migrazione verso ambienti Cloud e SaaS}
Per la migrazione verso il \textit{Cloud}, sono necessari \textit{tool} robusti in grado di lavorare sia con logiche \textit{batch} sia con \textit{API} e servizi \textit{web} (\textit{webservices} o \textit{WS} in breve),
adattandosi alle regole di \textit{business} e garantendo la qualità di ogni singolo dato prima ancora che venga scritto nel
database del nuovo \textit{tenant}.

